\chapter{Fitting Background}
\label{Chap:background}

The process of fitting the background in a diffraction pattern is potentially complex, in that GSAS-II provides four different computational mechanisms for background fitting and these are commonly used in combination. These will be described further, below. GSAS-II provides some unique tools can assist in background fitting. They are not needed for the simple cases, but may greatly simplify the problems associated with background modeling for the problem cases, where background is not a smooth slowly increasing or decreasing function.

The process of fitting a background is conceptually simple. Any scattering that is not Bragg scattering from the sample needs to be defined as part of the model used for fitting. The background model is added to the Bragg scattering to produce a simulation that reproduces the entire observed diffraction pattern. If the background is not well fit, it will degrade the overall fit and may cause parameters in the Bragg fitting to deviate from the correct values as the refinement attempts to minimize the differences between the observed and computed pattern. Most commonly, the parameters that would be affected by a poor background fit are the atomic displacement parameters ($\rm U_{iso}$) and/or peak shapes, but on occasion structural parameters can be biased. I should note that I am a firm believer that the final fits the background model should be optimized against the data rather than be fit by what the scientist thinks is correct. The best model predicted by the data is what is needed to minimize the overall fit while also obtaining chemically plausible results. While in the early stages of a fit, it may be best to fix the background to be in the locations that look right to our eye, my preference is that in the final stages, the profile be refined. There is one exception to this, which is when all the $\rm U_{iso}$ values refine to much lower than expected values. In this case it may be necessary to use a hybrid approach, as will be discussed in section \ref{negativeUiso} of this chapter.

Before discussing background fitting, I will define one term that is useful for these discussions, the baseline. I will speak of the baseline as function that describes the background as a function of TOF, $2\theta$, Q, etc. This is the average intensity well away from any Bragg peaks, but will be well below the data in the regions where peaks are found. Note that where noise levels are high, there may well be data points that are below the baseline, since the noise will scatter the data points around this average. Date may not return to the baseline for quite some distance from peaks, if the broadening has appreciable Lorentzian character (from sample broadening), since the Lorentzian line shape has very long peak tails. 

\section{Background Computation Modes}

The goal for a background fitting function is that it should produce a smoothly varying intensity curve over the range of $2\theta$ or TOF values and the complexity of the curve should depend on the number of terms for the 

GSAS-II provides several different background functions thapproaches for computing background scattering. 

\subsection{Smooth Refined Functions}

GSAS-II provides a number of functions that can be 

\subsection{Debye Scattering Peaks}

\subsection{Background Peaks}

\subsection{Fixed Background Histogram}


\section{Approaches to Fitting Background}

 Further, while sometimes all that is needed is to pick the right function and refine its coefficients, that does not always work, so GSAS-II provides both a manual and a semi-automated way to "bootstrap" the background. 

\subsection{Setting Data Limits}
One key step in background fitting is to use the right data limits for the refinement. It is common that the low Q portion of a diffraction pattern has a complex background. If there are no diffraction peaks in that region, one is only making the fitting process more difficult by including that portion of the pattern. Make sure to set the low Q data limit (low angle or long TOF time) so that the lowest-order observed diffraction peaks are included in the computation but do not include data at much lower Q than where the first peak reaches the baseline. While typically not a problem with respect to background fitting, including data at too high a Q value is also not a good idea. Look at the data at high Q and see if there are observable changes in intensity or if the intensities are too small to see relative to the noise level. Fitting noisy data can be helpful, but do not include data at high Q where there is only background to fit. The most important reason to not include this "flat"  data at high Q is not for background considerations, but due to the time needed to perform computations. As the range of the pattern increases in Q, the number of reflections increases by ~$\rm Q^3$ and the time needed to compute the pattern also increases on that order. This is most noticeable with TOF data, where a small change in the minimum TOF time can add a very large number of reflections to the pattern. I have seen least-squares refinement cycle times drop from minutes to seconds with the choice of sensible diffraction limits. 

Limits for the range of data to be used in fitting are stored on the Limits data tree entry and can be changed there, but can also be changed by dragging the green and red vertical dashed  lines from the plot for the background tree entry.

\subsection{Measuring Instrumental Background}

GSAS-II allows you to measure a diffraction pattern that will be subtracted from your observed pattern, point by point. Note that the two patterns must have the same point spacing and the same start and ending values. This can be a very useful thing, when working with "background-rich" data collection environments, common with \textit{in situ} experiments where ancillary equipment (furnaces, refrigerators/cryostats, reaction cells, etc.) can produce significant amounts of scattering that can have sharp features that are hard to reproduce. When I worked on PDF determination, where it was important to separate non-Bragg scattering from the sample from instrumental background, we collected the background scattering from the instrument without any samples and a pattern with the sample refrigerator but without any samples loaded, ideally with an empty sample container. One might think that simply subtracting the observed background that empty ancillary equipment should account for all background, but there are several issues that this will not address. The first is that the absorption and diffraction from the sample reduces the amount of radiation that can be scattered from the ancillary equipment. This means that the measured background will usually be a slight overestimate for what is actually observed from that source. In GSAS-II, one can specifiy 

\section{Negative $\rm U_{iso}$ values}\label{negativeUiso}
On occasion it is not possible to refine the peak shape, ADP’s (Uiso values) and background together. This is usually in materials that have many peaks, and where the peaks start to overlap so that beyond some value of Q there are no areas without reflections so the intensities are always above the baseline for a large region of the pattern. These areas without peaks define where the background must be located. This means that if the high-Q background were to be placed at a slightly lower level, a comparable fit can be obtained by increasing the intensity of the high-Q data (by lowering the Uiso values for the atoms) and adjusting the peak shapes to put more intensity into the peak tails. The symptom for this will be a refinement that seems quite good in quality, but has very low, often negative, U\textsubscript{iso} values. If a single U\textsubscript{iso} value is negative and all other atoms have expected values, this can mean that the atom type is wrong or the occupancy is too low and more scattering density is needed at that site, but when all U\textsubscript{iso} values are low this is an indication that there is not a single model to fit the data, but rather a range of values that produce similar fits. 

 My solution for this is to set all the U\textsubscript{iso} values to a value somewhat larger than I would expect based on my experience for the material type and data collection temperature and then fit the background and sample broadening. I then stop fitting the background but do fit the U\textsubscript{iso} values again. At this point the U\textsubscript{iso} values will typically refine to larger values. Do note in any publications that the background could not be refined in the presence of other parameters and needed to be fixed. 